\section{Classification of Signals and Systems}

\subsection{Signal Classification}

\subsubsection{Continuous-time and discrete-time signals}
\[
x(t), \qquad x[n]=x(nT_s), \qquad t\in\mathbb{R},\ n\in\mathbb{Z}.
\]
A continuous-time signal is defined for all real time values. A discrete-time signal is defined only at sample instants.

\subsubsection{Analog, sampled, quantized, and digital signals}
\[
\Delta V_{\mathrm{LSB}}=\frac{V_{\mathrm{FS}}}{2^N}.
\]
Analog signals are continuous in time and amplitude. Digital signals are discrete in time and amplitude, with one least significant bit set by the full-scale range and bit count.

\subsubsection{Even and odd signal components}
\[
x_e(t)=\frac{x(t)+x(-t)}{2}, \qquad
x_o(t)=\frac{x(t)-x(-t)}{2}, \qquad x(t)=x_e(t)+x_o(t).
\]
Even components satisfy \(x_e(t)=x_e(-t)\). Odd components satisfy \(x_o(t)=-x_o(-t)\).

\subsubsection{Periodic signal and fundamental period}
\[
x(t+T_0)=x(t), \qquad \omega_0=\frac{2\pi}{T_0}.
\]
The smallest positive \(T_0\) is the fundamental period. Periodic signals have discrete harmonic frequencies \(n\omega_0\).

\subsubsection{Signal time shift and time scaling}
\[
y(t)=x(t-t_0), \qquad z(t)=x(at).
\]
Positive \(t_0\) delays the signal. Time scaling by \(|a|>1\) compresses the signal, and \(a<0\) also reverses time.

\subsubsection{Unit step, ramp, and rectangular pulse}
\[
u(t)=
\begin{cases}
0, & t<0,\\
1, & t\geq 0,
\end{cases}
\qquad
r(t)=t\,u(t),
\qquad
p_T(t)=u(t)-u(t-T).
\]
The step gates causal signals. The ramp is the integral of the step, and a rectangular pulse is formed by subtracting delayed steps.

\subsubsection{Dirac impulse properties}
\[
\delta(t)=\delta(-t), \qquad
\int_{-\infty}^{\infty}\delta(t)\,dt=1, \qquad
\int_{-\infty}^{\infty}x(t)\delta(t-t_0)\,dt=x(t_0).
\]
The impulse is an even generalized signal with unit area. Its sampling property extracts the signal value at the impulse location.

\subsubsection{Impulse as derivative of step}
\[
\frac{d}{dt}u(t)=\delta(t), \qquad
\int_{-\infty}^{t}\delta(\tau)\,d\tau=u(t).
\]
Impulse and step are related by differentiation and integration. This relation is used when converting between impulse, step, and ramp responses.

\subsection{System Classification}

\subsubsection{Linearity of systems}
\[
T\{a x_1(t)+b x_2(t)\}=aT\{x_1(t)\}+bT\{x_2(t)\}.
\]
A linear system obeys superposition and homogeneity. Terms such as \(x^2(t)\), \(y^2(t)\), products of signals, and additive constants violate linearity.

\subsubsection{Time invariance}
\[
T\{x(t-t_0)\}=y(t-t_0) \quad \text{when} \quad T\{x(t)\}=y(t).
\]
A time-invariant system has unchanged behavior after a time shift. Explicit dependence on \(t\) in coefficients usually violates time invariance.

\subsubsection{Causality}
\[
y(t_0)\ \text{depends only on}\ x(t),\ t\leq t_0.
\]
A causal system does not require future input values. Terms such as \(x(t+t_0)\) with \(t_0>0\) are noncausal.

\subsubsection{LTIC differential equation}
\[
\sum_{k=0}^{N}a_k\frac{d^k y(t)}{dt^k}
=
\sum_{m=0}^{M}b_m\frac{d^m x(t)}{dt^m},
\qquad a_k,b_m=\text{real constants}.
\]
Linear time-invariant continuous-time systems in the course are represented by constant-coefficient differential equations.

\subsubsection{Filter type from low-frequency and high-frequency gain}
\[
H(0), \qquad \lim_{\omega\to\infty}H(j\omega).
\]
A lowpass filter has nonzero low-frequency gain and vanishing high-frequency gain. A highpass filter has vanishing low-frequency gain and nonzero high-frequency gain, while a bandpass filter vanishes at both extremes.
